% 原书第 557 页（PDF 第 570 页）
\BookPage{557}{570}
\BookHeading{乾隆及其扬州诗作}{Emperor Quianlong and His Poetry of Yangzhou\\[3pt]（乾隆及其扬州诗作）}{韩小敏*\\[2pt](Sherman Hsiao-min Han)}
\noindent Qing Gaozong (清高宗 A.D. 1711-1799) is probably one of the most prolific poets in the history of China. Starting in his youthful years until he passed away at the age of 89, the emperor presumably wrote more than 42,640 poems which were found in such special collections as The Collection of Le Shan Hall (乐善堂集), Ceremonious Records of the Imperial Tours to the South (南巡盛典), and Royal Productions of Gao Zong: Collection of Poetry and Essays of Ten Completions (高宗御制诗文十全集), plus the complete Royal Poetry Collections (御制诗) consisting of 498 volumes (Liu 253). Of all those documents, Ceremonious Records of the Imperial Tours to the South (南巡盛典) contains many poems directly related to the Prefecture of Yangzhou (扬州府) composed during his first four imperial tours taken place respectively in AD 1751, 1757, 1762, and 1765 (Sun 230). They can generally be divided into four major categories based on their themes: first, instructions to the local officials on the basic principles of governing, such as several "Verse to Instruct All Officials of High and Low Ranks in the South" (示江南大小诸吏韵Records 253 and 525), and "A Poem to Instruct All Involved Officials on My Order to Build Rock Embankment Outside of Gaoyou"(命筑高邮城外石隄诗示董事诸臣 Records 328); second, compositions directed to many local officials he met along the journeys regarding his personal feelings, such as "To Huang Tinggui, Governor General of Liang Jiang"(赐两江总督黄廷桂Records 116),"To Gao Jin, Governor General of Rivers and Canals of the South"(赐江南河道总督高晋Records 531), and "To Shen Deqian (赐沈德潜 Records 120); third, topics relevant to the livelihood of the people, such as "A Poem Written upon Giving Order to Waive Unpaid

\vfill
\begin{BookNotes}
\BookNote{}{作者任职于杨百翰大学夏威夷分校(Brigham Young University - Hawaii)。}
\end{BookNotes}

% 原书第 558 页（PDF 第 571 页）
\BookPage{558}{571}
\noindent Taxes from the Ripe Rice Farms in the South, and A Half of Taxes to All Prefectures and Counties in Jiangsu on the Course of My Tour"(降旨免江南熟田未完欠项并将江苏经过州县均免十分之五诗以志事Records 537),"Wood Dragon"(木龙Records 119), "Windmill"(风车 Records 122) and "Rain in the Evening: Two Poems"(夜雨二首Records 260); fourth, descriptions of the tours and landscapes of the south, especially those scenic places in Yangzhou (扬州), such as "Seeing the Relics in Weiyang"(维扬览古 Records 124), "Xiang Fu Temple" (香阜寺 Records124, 262, and 395), "Gao Min Temple"(高旻寺 Records 124, 186, and 265), "Tian Ning Temple"(天宁寺Records 125, 263, and 395), "Ta Wan Travel Palace"(塔湾行宫 Records 125 and 398), and "Ping Shan Hall"(平山堂Records 126, 263-64, 395-96, and 398), "Rainbow Bridge"(虹桥Records 126). This paper intends to analyze and evaluate the characteristics of each of those four categories; some of the representative works in each category will be translated into English based on their original poetic forms.

In the first category, the emperor meant to instruct his officials, through poetry, about his own beliefs in general or about his opinions of specific events probably because unlike edicts and official proclamations which involuntarily provoke resistance as a result of its bureaucratic nature, poetry invites its reader to think about, interpret, and understand the subject voluntarily and more in depth. The most prominent theme in the category is the emperor's "Verse to Instruct All Officials of High and Low Ranks in the South"(示江南大小诸吏韵Records 115, 253, 387, and 525) which was composed at the beginning of each and every of his first four tours to the South. In those works the emperor tells the local officials about the main purposes of his tours, which include to witness the lives of his subjects in the south and to personally inspect the preventive engineering works against the floods; in the mean time he also encourages those local officials to grasp his concerns for the people and to empathize with their pains and hardships. The next theme deals specifically with his instructions on constructions of the riverbanks, as seen in "A Poem to Instruct All Involved Officials on My Order to Build Rock Embankment Outside of Gaoyou"(命筑高邮城外石堤诗示董事诸臣


% 原书第 559 页（PDF 第 572 页）
\BookPage{559}{572}
\noindent Records 328) and "Instructions to Governor General Yin Jishan and Deputy Governor General Ji Huang While Viewing the River"(阅河示总督尹继善副总嵇璜Records 328).

The most popular theme in the first category attends to Qianlong's concise instructions to the various local officials showing them how to become a better administrator in their respective capacities. In his "Instructions to Jiangsu Education Administrator Zhuang Yougong"(示江苏学政庄有恭Records 121) the emperor teaches that the educators should "respect purity, truthfulness, elegance, and justice, /While reject the cattle devils and snake spirits"(所贵清真兼雅正，莫容牛鬼及蛇神); in his"Instructions to Governor General Huang Tinggui"(示总督黄廷桂 Records 117) he sets the guidelines for government aids to the flood victims: "it's better to over compensate the sufferers; /this policy you should pursue restlessly"(恤民宁过厚，咨尔推善行). By the same token, he explicates the essences of a fine government administrator in the following two poems:

"At Ta Wan Travel Palace the Third Time in Accordance with the Poem of the Royal Grandfather"

(塔湾行宫三依皇祖示江南吏韵)

\begin{BookQuote}
You officials of all ranks should learn to follow
\end{BookQuote}

\begin{BookQuote}
Those fundamental principles in governing:
\end{BookQuote}

\begin{BookQuote}
Show forgiveness by using the cattail whips;
\end{BookQuote}

\begin{BookQuote}
Have integrity similar to the bamboo pavilions.
\end{BookQuote}

\begin{BookQuote}
Always treat people's pain as if it's your own;
\end{BookQuote}

\begin{BookQuote}
Be constantly concerned about their grievances.
\end{BookQuote}

\begin{BookQuote}
Why devise various clever ways to collect taxes
\end{BookQuote}

\begin{BookQuote}
When half of the taxes levied have been waved.
\end{BookQuote}


% 原书第 560 页（PDF 第 573 页）
\BookPage{560}{573}
\begin{BookQuote}
载咨大小吏，为政有常经。恕不妨蒲示，节当如竹亭。莫忘侗痹念，恒且怨咨形。
\end{BookQuote}

\begin{BookQuote}
岂必催科巧，逋征半已停。(Records 539)
\end{BookQuote}

\begin{BookQuote}
"Instructions for Governor General Gao Bin While Viewing the River"
\end{BookQuote}

(阅河示总督高斌)

\begin{BookQuote}
People live prosperously in a nation rich in waters; The Golden embankment came from our ancestors.
\end{BookQuote}

\begin{BookQuote}
Coming to inspect, I expect you to do your duties well;
\end{BookQuote}

\begin{BookQuote}
The true value of achieving success lies in no actions.
\end{BookQuote}

\begin{BookQuote}
Uphold justice and make judgments on frivolous talks;
\end{BookQuote}

\begin{BookQuote}
Open your hearts to gather numerous new thoughts.
\end{BookQuote}

\begin{BookQuote}
You asked me what were my instructions for you?
\end{BookQuote}

\begin{BookQuote}
My ambition's shown in this five-character poem.
\end{BookQuote}

\begin{BookQuote}
泽国民生庶，金隄祖制垂。来巡期善守，底绩贵无为。持正裁浮议，开诚集众思。问予何训示，志在五言诗。(Records 119)
\end{BookQuote}

Emperor Qianlong made it a habit to always write a poem at wherever Emperor Kangxi had composed probably in order to show that he was a filial grandson who obediently following his late grandfather's footsteps. As a consequence, poems with the same title, such as "Ta Wan Travel Palace in Accordance with the Poem of the Royal Grandfather'(塔湾行宫依皇族祖诗韵) would appear in each and every of his inspection tours (Records 125, 264-66, 398-99, and 538-39). Nonetheless the emperor's instructions are very clear and forceful in both of these two poems. There are no flowery images or profound symbols. The language is plain and direct; its tone authoritative yet encouraging. In strict compliance with the tonal and rhyme scheme of a regulated


% 原书第 561 页（PDF 第 574 页）
\BookPage{561}{574}
\noindent verse, those two poems effectively conveys the message to the reader like most other works of this category. They may not necessarily be considered artistically impressive, yet they serve their purpose quite remarkably. The second category of Emperor Qianlong's poetry of inspection tours consists of a great number of works bestowed upon those current or retired officials--some of whom lived or stationed in the region, some came to meet with him from other regions at the time. These works are different from those of the previous category in that the emperor treats their receivers more as a literary partner than as a subject who pledges utter obedience to the throne. A popular theme in the category is the emperor's greetings to his officials. For instance, he dedicated several poems to Shen Deqian (沈德潜), a retired official and literary partner upon whom the emperor repeatedly bestowed rewards and honors during all of his four inspection tours (Sun 229-30), to reminiscence of the past and celebrate his longevity (Records 120, 262, 548); sometimes he ordered Shen to reciprocate with the same topics and equal number of works, such as the twelve poems composed upon his journey to She Mountains (摄Records 603). Another popular theme in the category is the emperor's congratulatory notes to his high ranking court officials upon their promotions, such as his "To Yin Jishan, Governor General of the South and Grand Secretary"(赐江南总督大学士尹继善Records 527) for his being promoted to the post of grand secretary, and "To Assistant Grand Secretary, Minister of Justice, also the Governor of Jiangsu, Zhuang Yougong"(赐协办大学士刑部尚书仍管江苏巡抚庄有恭Records 527) for his additional title of assistant grand secretary. However, the majority of poems in the second category are given to the former and current officials praising their job performances. They include, among others, "To Gao Jin, Governor General of the Rivers and Canals of the South"(赐江南河道总督高晋Records 389, 531), "To Bai Zhongshan, Governor General of the Rivers and Canals" (赐河道总督白钟山Records 331), "To Chen Hongmou, Governor of Suzhou"(赐苏州巡抚


% 原书第 562 页（PDF 第 575 页）
\BookPage{562}{575}
\noindent 陈宏谋 Records 395),"To Li Yinpei, Education Administrator of Jiangsu "(赐江苏学政李因培Records 274),"To Liu Yong, Education Administrator of Jiangsu"(赐江苏学政刘墉 Records 395), and "To Huang Tinggui, Governor General of Liang Jiang"(赐两江总督黄廷桂Records 116) translated as follows:

\begin{BookQuote}
Traffics were cleared to welcome the imperial carriage;
\end{BookQuote}

\begin{BookQuote}
The general with a royal tally has long been on his tours.
\end{BookQuote}

\begin{BookQuote}
Knowing my diligent care for the welfares of the people,
\end{BookQuote}

\begin{BookQuote}
He's led his worthy officers improving the villagers' lives.
\end{BookQuote}

\begin{BookQuote}
What is the need to set up all those colorful silk curtains?
\end{BookQuote}

\begin{BookQuote}
I only enjoy looking at the sights of the scaly rice fields:
\end{BookQuote}

\begin{BookQuote}
Carefully experiencing the hard and soft local customs,
\end{BookQuote}

\begin{BookQuote}
And inquiring fully about the crops abundant and lean?
\end{BookQuote}

\begin{BookQuote}
Perceiving to stay with golden mean will result in success;
\end{BookQuote}

\begin{BookQuote}
Prejudice will disappear naturally when ego is evacuated.
\end{BookQuote}

\begin{BookQuote}
The residents and cultures are prosperous in this place;
\end{BookQuote}

\begin{BookQuote}
Encourage them to restore to their original pureness.
\end{BookQuote}

\begin{BookQuote}
迎銮卫警跸，建节久旬宣。体我勤民德，嘉乡率属贤。何须张锦帐，惟喜阅鳞田。
\end{BookQuote}

\begin{BookQuote}
细验刚柔俗，周咨丰歉年。咸中方有庆，虚己自无偏。此地人文盛，还淳尚勉旃。
\end{BookQuote}

The use of diction, imagery, and form as shown in this poem is not much different from that of the works in the first category. The language in the third and sixth couplets, for instance, is simple and conversational to some degree. The images of "the imperial carriage," "colorful curtains,"and "rice fields" are certainly nothing striking.


% 原书第 563 页（PDF 第 576 页）
\BookPage{563}{576}
\noindent The major difference between the works of those categories lies in the poetic tones. Unlike the formal and sometimes austere tones shown in the last group, the tones of this group are relatively informal, relaxing, and friendly, thus convey a feeling of closeness between the writer and his reader. As a result, those works serve the emperor well as a literary means to symbolically endear himself to his officers while encouraging them to further improve their administrative performances.

The third category consists of poems that mainly show the emperor's concerns for the welfares of his subjects. The most notable theme in the category is the close attentions he paid to the rainfalls during his inspection tours. Assuming that rain waters would help the crops and thus improve people's lives, he would frequently write a poem expressing his happiness when it rained in a place he happened to be visiting at the time. In his"Two Poems about Rains at Night" (夜雨二首 Records 260) composed near Ping Qiao (平桥), the emperor says, "This rain brings thorough moist; /It'll comfort the old farmers "(此雨设遍沾，可以慰农父),and "Myriads of silky rains in February/Would benefit farmers everywhere"(二月雨霏丝，到处应益农). Further still, while on the inspection tours he would write poems about rainfalls in the other regions or even at the capital whenever it was reported to him. Those include "Remark Happily of Yin jishang and Chen Hongmou Reporting Suzhou Receives Rains"(尹继善陈宏谋奏报苏州得雨志喜 Records 445), "Governor Aertai Reporting Shandong Receives Rains"(山东巡抚阿尔泰奏报得雨Records 463),"Governor Cui Yingjie Reporting Shandong Receives Rains"(山东巡抚崔应阶奏报得雨 Records 533),"Governor Hu Baoquan Reporting Henan Receives Rains"(河南巡抚胡宝泉奏报得雨Records 465),"Heard That the Capital Receives Rains"(闻京师得雨 Records 531), "Again Heard That the Capital Receives Rains"(复闻京师得雨Records 286, 315, 609) and "Remark of Thoughts about the Capital Receives Rains"(京师得雨志怀 Records 341).

The next popular theme in the category comprises of farming and the machinery that is closely involved with lives of farmers and other commoners in the south. In "Wheat


% 原书第 564 页（PDF 第 577 页）
\BookPage{564}{577}
\noindent Flowering"(麦秀 Records185) he writes:

\begin{BookQuote}
Several fine and sunny days hasten wheat flowering;
\end{BookQuote}

\begin{BookQuote}
This proves that the seasons come early in the south.
\end{BookQuote}

\begin{BookQuote}
Underneath straw eves, convinced of harvest coming,
\end{BookQuote}

\begin{BookQuote}
I know, in distance, the price in the market will not rise.
\end{BookQuote}

\begin{BookQuote}
数日好晴催麦秀，南方候早信今征。茅檐自觉收堪待，蒸肆遥知价不增
\end{BookQuote}

Similar lines appear in "The Wheat Seedling"(麦苗254) and "The Wheat Colors"(麦色 Records 391), such as "Seeing the wheat colors while crossing the river /I realize the climate difference in the North and the South"(过河逢麦色，寒暖北南殊). In"Windmill" (风车Records 122) he explains in details how the windmill operates even better than the hydraulic mills for irrigations; in "Wood Dragon"(木龙 Records 119) and two other poems, "Composed to Instruct Gao Jin, Administrator General of the Rivers"(木龙示总河高普作 Records 391, 530), the emperor describes how "the wood dragon," which was first invented by Chen Yaozuo (陈尧佐) in the Song dynasty (宋), should work in clearing the accumulated sand in Tao zhuang(陶庄).

The most interesting theme in this category is the emperor's self awareness of how his decisions on the inspection tours could bring direct impacts upon his subjects as shown in the following works:

\begin{BookQuote}
"A Poem to Remark My Giving Order to Waive Unpaid Taxes from the Ripe Rice Farms in the South, and A Half of Taxes to All Prefectures
\end{BookQuote}

and Counties in Jiangsu on the Course of My Tour"

\begin{BookQuote}
(降旨免江南熟田未完欠并将江苏经过州县均免十分之五诗以志事)
\end{BookQuote}


% 原书第 565 页（PDF 第 578 页）
\BookPage{565}{578}
\noindent A total of one point three million of taxes were waved before departure;

Those were referred to the remaining taxes owed by the flood victims.

Once in the region I found out that more of the waivers went to the North;

Only around thirty thousand of taxes had actually been waived in the South.

In this case, the North of the River received the major portion of benefits;

The South received fewer benefits because over there no floods took place.

I asked whether there was any tax owed from those rice fields that were ripe,

And was told that the total owed from those fields was merely forty thousand.

It seems the benefits to them would still be light even all of that is dropped;

So I decided to benefit everyone equally by waiving one half of all their taxes.

Confucius once said not to give wealth to those who are already wealthy;

To benefit the others in that manner is not behaving like a man of pure heart.

To believe in storing vast wealth in the common households is perfectly fine;

Who says that my decision was only meant to gratify and please the people?

启銮前免百卅万，实指灾田所欠尾。入疆询知江北多，江南才三万而已。

是则江北被恩重，江南无灾恩薄矣。因问熟田或有逋，亦不过四万斯止。

槩与蠲之仍觉轻，十分免五惠同视。宣尼曾云不继富，益于他人非赤子。

藏富闾阎信亦佳，讵曰驩虞博民喜。(Records 537)


% 原书第 566 页（PDF 第 579 页）
\BookPage{566}{579}
"Random Talks"(杂言)

Yesterday blew the southeastern winds;

The boat seemed to sail rather slowly.

Today came the northeastern winds;

The boat seems to sail exceedingly fast.

Presently the wind was in its full force;

The boast now sailed like it was flying.

It charged ahead indifferently in the middle of the river;

Ahead of it for about a few miles sailed a small boat.

In an instant it had already been reached by our boat;

The panicky small boat tried to dodge us desperately.

If unable to dodge in time, it would surely be crushed;

Seeing all this alone by the railing, my heart was startled.

Thus is the disparity between the large and small forces;

What at an emperor can enforce, nothing could compete.昨日东南风，舟行诚觉迟。今日东北风，舟行迅似驰。。。须臾风力足，越驰乃如飞。大船中流若无事，小船前行数里计。瞥眼大船已及之，音后小船欲回避。回避不能将压临，凭览独觉惊吾心。大小不伦势若此，可恃为君莫吾比。

(Records 534)


% 原书第 567 页（PDF 第 580 页）
\BookPage{567}{580}
\noindent In the former poem the emperor tries to defend his order of waiving taxes to the farmers in Jiangsu on the basis of Confucius' teaching and the economic theory of conserving riches in the people instead of the government. The Chinese characters of"not a man of pure heart" (非赤子) in the sixth couplet could refer to the impropriate deed of Gongxichi (公西赤) who was the target of Confucius' rebuke of trying to gain more stipends for his already well paid position; those of "huan yu" (萝虞) might be a pun with the historical character of Feng Xuan (冯驩) in the Warring States (475-221 BC), who rightly forfeited the interests and principals owed to Meng chang jun (孟尝君) by the poor of Xue (薛), and thus made Meng a well liked and respected benefactor upon whose return to the area after a forced retirement from the premiership of Qi (齐). In the latter work, the poet cautions himself to exercise his power responsibly as he learns the imbalance between the powers of an emperor and those of the civilians by witnessing the inevitable crash between the large royal boat he was boarding and a small sail boat ahead of it. In spite of the unusually easy language, these two poems demonstrate relatively deep thoughts and witty references not commonly seen in his other poems of Yangzhou.

Included in the last category are poems depicting the historical and scenic poínts of Yangzhou and its neighboring towns. According to Revised Gazetteer of Yangzhou Prefecture (重修扬州府志"Ink Marks of Kangxi and Qianlong"), Qianlong wrote a total of 115 poems on the scenic temples in Yangzhou alone, including 16 poems on Tian Ning Temple (天宁寺), 22 on Gao Ming Temple (高旻寺), 57 on Ping Shan Tang Temple (平山堂),four on Guanyin Shan (观音山), four on Fa Hai Temple (法海寺), one on Chong Ning Temple (重宁寺),three on Chan Zhi Temple (禅智寺), four on Xiang Fu Temple (香阜寺), and four on Hui Yin Temple (慧因寺). An additional number on Gaoyou Lake (高邮湖), Shao Bo Lake (邵伯湖 123-25, 262), Lu Jin Temple (露筋祠 Records 261), Wenyou Tai (文游台 261), Ta Wan Travel Palace (塔湾行宫Records 125, 264-66, 398, 538), "Rainbow Bridge" (虹桥 Records 126), and "Lian Xing Temple" (莲性寺 Records 263) are found in Ceremonious Records of the Imperial Tours to


% 原书第 568 页（PDF 第 581 页）
\BookPage{568}{581}
\noindent the South(南巡盛典) and some local gazetteers (Gazetteer of Shao Bo Town邵伯镇志461-63).

Despite the large number, they can again be generally divided into two groups based on the poetic themes: one consists of poems that purely focus on the descriptions of the scenic boat tours and the spots that the emperor personally visited, such as the following works:

"Seeing the Relics in Wei Yang"(维扬览古)

\begin{BookQuote}
Untied and departed from Qin You in the early morning,
\end{BookQuote}

\begin{BookQuote}
The royal boat stations at Wei Yang under the setting sun.
\end{BookQuote}

\begin{BookQuote}
The hills with spears and banners look like the Shu region;
\end{BookQuote}

\begin{BookQuote}
Those mounds on the mountain seem like they're afloat.
\end{BookQuote}

\begin{BookQuote}
Commander Xie's calmness is still being talked about;
\end{BookQuote}

\begin{BookQuote}
Governor Ou's works further remind us of their writer.
\end{BookQuote}

\begin{BookQuote}
There're smoky willows on the lonely Sui embankment;
\end{BookQuote}

\begin{BookQuote}
Who saw the tumultuous fireflies at this autumn night?
\end{BookQuote}

\begin{BookQuote}
清晨解缆发秦邮，落照维扬驻御舟。冈有枪旗宛是蜀，山才培壤亦称浮。
\end{BookQuote}

\begin{BookQuote}
建牙镇静犹传谢，通守文章更纪欧。寂寞隋隄烟柳在，秋宵谁见乱萤流。
\end{BookQuote}

(Records 124)

"Xiang Fu Temple"(香阜寺)

\begin{BookQuote}
Passengers disembarked from a boat
\end{BookQuote}

\begin{BookQuote}
Pulled by the ropes in the late spring.
\end{BookQuote}


% 原书第 569 页（PDF 第 582 页）
\BookPage{569}{582}
\noindent After several legs of water county,

It's a pleasure to see the mountains.

Though nothing's unusual about them,

Their name's allied with fresh coolness.

Writings are not labeled as large or small;

Only people are obsessed with positions.

Temporarily I've lost my thoughts in a room

Where spikes of smoke decorate the incenses.

维舟对晚春，登陆临初地。水乡经数程，见山已快意。土阜纵无奇，名与清凉配。

文殊无大小，而人纷位置。丈室暂僵徨，炉香缀烟穗。(Records 124)

"Gao Ming Temple"(高旻寺)

A green lotus temple sits in the tranquil woods;

A Buddhist tower stands in midst of the blue sky.

The famous bay's truly worthy of its reputations;

The story of goose burial deserves to be retold.

Unable to look far as one would like in a clear day,

Who'd forget to explicate it with words and images?

The Golden Mountain, being an uninvited guest

Could only hide temporarily behind the river smokes.


% 原书第 570 页（PDF 第 583 页）
\BookPage{570}{583}
\begin{BookQuote}
兰若青莲宇，浮图碧落天。名湾真不愧，埋雁亦堪传。未纵晴明望，谁忘言象诠。
\end{BookQuote}

\begin{BookQuote}
金山不速客，暂尔隐江烟。(Records 124)
\end{BookQuote}

"Rainbow Bridge"(虹桥)

\begin{BookQuote}
Spring water with green waves is drinking the long rainbow;
\end{BookQuote}

\begin{BookQuote}
Slowly pulled by the silk cables, the boat sails like in a mirror.
\end{BookQuote}

\begin{BookQuote}
It's as if we are really moving across in a horizontal picture;
\end{BookQuote}

\begin{BookQuote}
Blowing in my face are fragrant breezes from Mountain Ping.绿波春水饮长虹，锦缆徐牵碧镜中。真在横披图里过，平山迎面送香风。
\end{BookQuote}

(Records 126)

"Shao Bo Town"(邵伯)

\begin{BookQuote}
Tai Fu embankment sits by the riversides of the green waters;
\end{BookQuote}

\begin{BookQuote}
It greatly benefits residents of Shao Bo who sing praise of Tang Yin.
\end{BookQuote}

\begin{BookQuote}
As the orchid boat passes slowly by, I think of those old footsteps,
\end{BookQuote}

\begin{BookQuote}
And hear the exquisite music of strings and bamboo at Dong Shan.
\end{BookQuote}

\begin{BookQuote}
太傅隄存绿水浔，惠方邵伯颂棠阴。兰舟缓过思遗躅，觉胜东山丝竹音。
\end{BookQuote}

(Records 124-25)

Probably due to the nature of the subjects, the overall dictions used in the works of this category are more suggestive and graceful than those of the other three. Though not necessarily outstanding while compared to those of the masterpieces, the language in those poems does appear to be adequate, natural, and eloquent; and they do occasionally provide some impressive imagery, such as "the smoky willows on the


% 原书第 571 页（PDF 第 584 页）
\BookPage{571}{584}
\noindent lonely embankment (寂寞隋隄烟柳在)"and "the tumultuous fireflies at an autumn night (秋宵谁见乱萤流)"in "Seeing the Relics in Wei Yang "(维扬览古), "spikes of smoke decorating the incenses (炉香缀烟穗)"in "Xiang Fu Temple"(香阜寺), and"spring water with green waves drinking the long rainbow (绿波春水饮长虹)" in"Rainbow Bridge"(虹桥).

The other group consists of works that show the emperor's political thoughts in addition to the usual description of landscapes as manifested in the following examples:

\begin{BookQuote}
"Poems Composed Contemporaneously upon Changing Boat
\end{BookQuote}

from Gao Bridge to Tian Ning Temple"(自高桥易舟至天宁寺行馆即景咏)

\begin{BookQuote}
Ancient folks spoke of the March flowers in the mists;
\end{BookQuote}

\begin{BookQuote}
Yangzhou has had many strings and pipes ever since.
\end{BookQuote}

\begin{BookQuote}
I'd like to ban them so people could return to being pure,
\end{BookQuote}

\begin{BookQuote}
Yet fear this could hurt swarms of workers in the business.
\end{BookQuote}

\begin{BookQuote}
三月烟花古所云，扬州自昔管弦纷。还淳拟欲申明禁，虑碍翻殃谋食群。
\end{BookQuote}

\begin{BookQuote}
Those welcoming programs on both riversides appeal to me not,
\end{BookQuote}

\begin{BookQuote}
Yet it's hard to reject the eagerness and hospitality of the crowds.
\end{BookQuote}

\begin{BookQuote}
I was nonetheless very pleased to pass by the picture of plowing
\end{BookQuote}

\begin{BookQuote}
And weaving since it's crucially related to people's clothes and foods.
\end{BookQuote}

\begin{BookQuote}
夹岸排当实厌闹，殷勤难却众诚殚。却从耕织图前过，衣食攸关为喜看。
\end{BookQuote}

(Records 535)


% 原书第 572 页（PDF 第 585 页）
\BookPage{572}{585}
"Tian Ning Temple"(天宁寺)

Tian Ning Temple outside of Tian Ning Gate is

The oldest and best flower palace in Guang Ling.

Moon-pointing branch's seen a hundred generations;

There rises a three-story tower that houses the sutras.

To endear myself to the people I suddenly think of Su Shi;

To pass on the doctrines need we always ask for Hui Neng?

I have had many worries, and cared not being happy alone;

How'd I indulge myself in poetry in this sphere of incenses?

天宁门外天宁寺，最古花宫冠广陵。指月禅枝阅百代，度经阿阁耸三层。

亲民乍为思苏轼，传法何须问慧能。我自多忧非独乐，岂能香界镇吟凭。

(Records 262)

"Ping Shang Hall"(平山堂)

Pines and bamboos grow thickly at the west side of West Temple

Where I toured, trying to find the old vestiges of poet Ouyang.

The Southern mountains and landscapes display boundless grace;

The springtime of the second month yields unrestricted beauty.

Hung underneath the eaves are four words by my grandfather,

Which serve as the maxim for all officials for a thousand years.

Looking up to his grand examples in everywhere on this tour,


% 原书第 573 页（PDF 第 586 页）
\BookPage{573}{586}
\begin{BookQuote}
Here I especially would like to appeal to the outstanding hearts.
\end{BookQuote}

\begin{BookQuote}
西寺西头松竹深，欧阳旧迹试游寻。江南山色秀无尽，二月韶光美不禁。
\end{BookQuote}

\begin{BookQuote}
四字檐端垂圣藻，千秋座右揭官箴。春巡处处前巍仰，到此尤廑吁俊心。
\end{BookQuote}

(Records 264)

The emperor explicitly tries to make his concerns for the livelihood of his subjects as an acceptable excuse for the exorbitant spending on account of his inspection tours in the first two poems. The same kind of guilt feelings probably made him reminiscence Su Shi's example of endearing himself to the people, and state that "I care not being happy alone(我自多忧非独乐)"in "Tian Ning Temple"(天宁寺). The political consideration in the last poem is slightly different from that of the other three. Here the emperor probably intends to make his extreme respect for his late grandfather, like his poetry at Ta Wan Travel Palace discussed earlier, as an example of the Confucian teaching of filial piety for scholars as well as commoners.

In conclusion, many critics suggest that in spite of the large volumes, very few of Emperor Qianlong's poetic works are of high quality (Sun 253), and even though considered a very well read and brilliant scholar, he is hardly regarded as a talented poet (Liu 371). However, based on the thematic studies of his poetry of Yangzhou, the emperor has been very skillful in making poetry as an effective means to convey his message, both explicitly and implicitly, to his readers: the instruction poetry not only give clear directions to the court officials, but also leave room for them to further interpret its deeper meanings as a poem. The works in the second category delicately endear the emperor to the local receivers of current and retired officials, and thus makes him appeal to the businessmen and scholars in the South, a main controversial target in the early stage of his rule (Bai 130). The third category unabashedly informs the people of his deep concerns for and undivided dedications to their living conditions and welfares. Similar devotions and moral teachings are seen even in his


% 原书第 574 页（PDF 第 587 页）
\BookPage{574}{587}
\noindent works of sightseeing and leisure in the last category. The emperor once stated that "I have never cared about the literary clichés of winds, clouds, moon, and dews while writing poetry, yet there will always a poem to remark the important political events or ceremonies (予向来吟诗，不屑为风云月露之辞。每有关政典之大者，必有诗记事Liu 253-54). Consequently, Qianlong's works do not necessarily earn him the title of a great poet, but they do successfully create an image of a good, which was probably what he had intended in the first place.


% 原书第 575 页（PDF 第 588 页）
\BookPage{575}{588}
Works Cited

Akedang, Axiu (阿克当阿修) and Yao Wentian (姚文田). Gazetteer of Yangzhou Prefecture: Jiaqing Revised Edition (嘉庆重修扬州府志). 2 vol. Yangzhou: Guangling (广陵),2006.

Bai, Xinliang (白新良). The Biography of Qianlong (乾隆传). Sheyang (沈阳): Liaoning Education Publishers(辽宁教育出版社),1990.

Gao, Dehai (高德海) et al. Gazetteer of Shao Bo Town (邵伯镇志). Nanjing: Jiangsu People's Publishers (江苏人民出版社),1996.

Gao, Jin (高普). Ed. Ceremonious Records of the Imperial Tours to the South (南巡盛典). 4 vol. Taipei: Wen Hai (文海), 1966.

"How Much Ink Mark of Kangxi and Qianlong in Yangzhou?"("康乾在扬墨迹知多少?"). Yangzhou Evening News (扬州晚报).24 November 2007.

Liu, Zhengguo (刘正国). The Life Philosophy of Qianlong (乾隆的人生哲学). Taipei: Han Yang (汉扬), 1997.

Sun, Wenliang (孙文良). Emperor Qianlong (乾隆皇帝). Taipei: Zhi Shu Fang (知书房), 2001.

Xiong, Hui (恢). Ed. Royal Productions of Gao Zong: Collection of Poetry and Essays of Ten Completions (高宗御制诗文十全集). Taipei: He Ji Shi Lin (合记士林), 1962.


% 原书第 576 页（PDF 第 589 页）
\BookPage{576}{589}
